\subsection{Experiment: how much does the prompt contract matter?}
\label{sec:prompt}
The failure analysis suggests that many errors arise because the original extraction prompt does not tell the model how to express things the IR \emph{can} express: it lists node tags without the time and space coordinates, vector construction, initial conditions or the form of a constitutive relation, and its only format example is a scalar Poisson model. To test this hypothesis we extended the system prompt by a short block (Appendix~\ref{app:supp}; about 350 tokens) that documents exactly these node types and rules, and re-ran the six documents in HTML for the four models that finished in the original experiment (the runaway-reasoning model was not repeated). Nothing else changed: same documents, same temperature, same scorer.

\begin{table}[t]
\centering\small
\caption{Effect of the extended prompt contract (HTML documents, six per model): documents valid, correct and ran (solver compiled and ran) with the original and with the extended prompt.}
\label{tab:v2}
\begin{tabular}{@{}lcccccc@{}}\toprule
& \multicolumn{3}{c}{original prompt} & \multicolumn{3}{c}{extended prompt} \\ \cmidrule(lr){2-4}\cmidrule(lr){5-7}
LLM & valid & correct & ran & valid & correct & ran \\ \midrule
Llama 3 8B & 3 & 2 & 2 & 4/6 & 0/6 & 0/6 \\
Qwen3-Coder 30B & 2 & 1 & 1 & 5/6 & 4/6 & 3/6 \\
Devstral Small 2 24B & 3 & 1 & 2 & 4/6 & 4/6 & 3/6 \\
Qwen3.6 27B & 6 & 4 & 3 & 6/6 & 6/6 & 4/6 \\
\bottomrule\end{tabular}
\end{table}

\begin{table}[t]
\centering\small
\caption{Per-document outcome with the extended prompt (\cmark\ correct, $\sim$ valid but not equivalent, \xmark\ rejected at the trust boundary).}
\label{tab:v2matrix}
\begin{tabular}{@{}lcccccc@{}}\toprule
LLM & Heat (NAFEMS) & Poisson & Transient heat & Elasticity & Stokes & Allen--Cahn \\ \midrule
Llama 3 8B & $\sim$ & $\sim$ & $\sim$ & \xmark & $\sim$ & \xmark \\
Qwen3-Coder 30B & \cmark & \cmark & \cmark & \cmark & \xmark & $\sim$ \\
Devstral Small 2 24B & \cmark & \cmark & \cmark & \xmark & \cmark & \xmark \\
Qwen3.6 27B & \cmark & \cmark & \cmark & \cmark & \cmark & \cmark \\
\bottomrule\end{tabular}
\end{table}

Across the four models, the number of documents (out of 24) that were valid rose from 14 to 19, the number of correct ones from 8 to 14 (from 33\% to 58\%), and the number whose solver compiled and ran from 8 to 10 (Table~\ref{tab:v2}). The effect is concentrated where the original prompt had a gap (Table~\ref{tab:v2matrix}). The transient-heat document, solved by none of the models before, was now correct for three of four, and the three corresponding solvers ran and reproduced the gold output exactly: the models used the time coordinate. For Allen--Cahn the 27B model now extracted the initial field and the relation, and its solver output equals the gold output to $4\cdot10^{-17}$. Qwen3-Coder and the 27B model now produced a correct elasticity model. The effect is not uniform: the smallest model got \emph{worse} (from 2 to 0 correct): with a prompt about twice as long it produced equations that fail the scorer and, in two runs, programs whose generated C++ did not compile, so a longer contract is not automatically better for an 8B model.

\paragraph{The contract has to match the backend.} Only 10 of the 14 correct extractions ran. The reason is instructive: our own extension told the models to write a zero vector as a \code{VectorExpr} of zeros, which the generator rejects (it accepts a \code{Constant} $0$ for the zero vector), and several models declared the pressure in $L^2$, which the generator does not support. Two of the four correct elasticity and Stokes extractions were therefore refused by the \emph{backend}, i.e., the prompt promised more than the generator delivered. We corrected both points (a ``v3'' prompt: zero vector as \code{Constant} $0$; all fields declared \code{H1}, the spaces being chosen by the solver) and re-ran the two vector-valued documents with three models (Table~\ref{tab:v3}). Both Stokes extractions that were correct now ran (they are correct only modulo the incompressibility constraint, with a 16\% velocity and a 33\% pressure difference to the gold discretization, the same kind of difference as in Section~\ref{sec:e2e}); the 27B model's elasticity extraction was correct but the generated program did not compile (the generator emitted invalid C++ for that equivalent formulation), and Qwen3-Coder's Stokes extraction was accepted by every check, ran, and solved a different problem (a pressure that differs from the gold one by more than 100\%), a second silently wrong run in the whole study. The lesson is that the prompt contract, the IR schema and the generator's actual scope should be derived from one specification; today they are three hand-maintained texts.

\begin{table}[t]
\centering\small
\caption{Second revision of the prompt (v3; elasticity and Stokes only, HTML). \cmark\ correct, $\sim$ valid but not equivalent, \xmark\ rejected; ``ran'' and ``exact'' count documents whose solver compiled and ran, and whose output equals the gold output.}
\label{tab:v3}
\begin{tabular}{@{}lcccc@{}}\toprule
LLM & Elasticity & Stokes & ran (of 2) & exact (of 2) \\ \midrule
Qwen3-Coder 30B & $\sim$ & $\sim$ & 1 & 0 \\
Devstral Small 2 24B & \xmark & \cmark & 1 & 0 \\
Qwen3.6 27B (thinking) & \cmark & \cmark & 1 & 0 \\
\bottomrule\end{tabular}
\end{table}

\paragraph{Caveat.} The extended prompts were written \emph{after} the original failures were analysed and evaluated on the same documents. The improvement shows that the prompt contract is a first-order lever, and it is probably an upper bound for what a generic revision would achieve on new documents. A fair evaluation requires a held-out corpus (Section~\ref{sec:outlook}).

